\section{Дизайн исследования и заморозка протокола}
\label{sec:design}

Этот раздел намеренно стоит перед результатами. Дизайн был зафиксирован
письменно до начала серии, и именно последовательность
\emph{гипотеза $\rightarrow$ замороженный метод $\rightarrow$ данные
$\rightarrow$ интерпретация} даёт читателю возможность отличить находку от
истории, собранной задним числом.

\subsection{Что было заморожено и когда}

Предрегистрация — артефакт репозитория, датированный раньше первого запуска
серии\code. В ней зафиксированы вопрос, гипотезы с условиями их
фальсификации, вмешательство, зависимые переменные, вмешивающиеся факторы и
контроль над ними, объём выборки с обоснованием мощности, правило остановки,
статистические тесты и политика исключений. Раздел поправок в ней пуст: после
первого запуска ничего не менялось.

Реальную работу делают три обязательства:

\begin{enumerate}
\item \textbf{Вердикт внешний.} Исход запуска определяет код завершения
      собственного верификатора задачи, а не отчёт агента о своём успехе.
      Запуск без вердикта верификатора в набор данных вообще не попадает.
\item \textbf{Тесты выбраны раньше данных.} Точный парный тест Макнемара — для
      способности решать задачи, точный знаковый тест на попарных отношениях —
      для стоимости, ранговая корреляция Спирмена — для тренда по горизонту.
      Ни один тест не подбирался задним числом, после того как эффект уже был
      замечен.
\item \textbf{Исключения определены до того, как увидели подсчёты.}
      Правило контаминации (\S\ref{sec:pitfalls}) и правило схлопывания
      дубликатов были написаны заранее, до применения. Второе оставляет в
      выборке \emph{самый ранний} запуск повторно запущенной ячейки, так что
      повтор не может протолкнуть себя в выборку победой.
\end{enumerate}

\subsection{Гипотезы и условия их опровержения}

\begin{description}[style=nextline]
\item[H1 (стоимость).] Медианное отношение стоимости только-состояния к
  контрольной группе превышает 1.0. \emph{Опровергается}, если отношение
  $\le 1.0$, знаковый тест не достигает значимости, или направление
  расходится между тремя испытаниями.
\item[H2 (горизонт).] С ростом горизонта отношение стоимости падает и
  пересекает 1.0: накладные расходы амортизируются на длинных задачах.
  \emph{Опровергается}, если корреляции нет, она положительна, или ни разу
  не приближается к 1.0 в наблюдаемом диапазоне.
\item[H3 (способность).] Приводится с доверительным интервалом; считается
  недостаточно мощной, если наблюдаемое несовпадение не превышает
  предвычисленное требование в шесть односторонних несовпадающих пар.
\end{description}

H2 — гипотеза, которую это исследование опровергает предрегистрированным
образом; мы сообщаем об этом прямо в \S\ref{sec:economics}, а не отодвигаем
в сторону.

\subsection{Множественные сравнения}

Предрегистрированы были три гипотезы, и именно они несут подтверждающий вес
исследования. Всё остальное, для чего статья приводит $p$-значение —
разбивки по стратам, ступени лестницы \S\ref{sec:composition}, сравнения
стоимости по условиям, — вторично и разведочно, и поправка на семейственность
сравнений к этому не применяется.

Мы проговариваем это явно, потому что \S\ref{sec:gaps} критикует оригинал
именно за отсутствие такой оговорки, а стандарт должен быть один для обеих
статей. Следствие конкретно: ключевой результат переживает любую поправку,
которую к нему можно применить ($p={}$\McnPTailZero{} по основному сравнению
против порога Бонферрони $0.05/3$ для подтверждающего семейства), а вот
несколько вторичных результатов лежат достаточно близко к $0.05$, чтобы её
не пережить. Там, где вторичный результат несёт вес для утверждения статьи,
мы говорим об этом прямо и трактуем его как предварительный, а не
установленный.

\subsection{Заранее объявленная ловушка артефакта}

Одна ловушка была названа заранее именно потому, что она грозила
приукрасить результат: \emph{более низкая} стоимость на длинных задачах в
руке только-состояния могла бы означать не амортизацию, а всего лишь более
частое исчерпание бюджета ходов — то есть цензурирование. Поэтому
предрегистрация исключает запуски, остановленные по бюджету, из агрегации
стоимости и приводит их число по каждой страте отдельно. Ловушка оказалась
реальной: распределение причин остановки в Таблице~\ref{tab:stop} показывает
остановки по бюджету, сконцентрированные именно в этой руке.

\subsection{Охват}

Эта серия отвечает на вопросы о стоимости и горизонте и приводит способность
решать задачи. Чувствительность к схеме и композиция состояния — отдельные
дизайны с собственными предрегистрациями; их статус дан в
\S\ref{sec:composition}.
\pending{если владелец решит опубликовать только E1, этот абзац станет
опережающей ссылкой на вторую статью}
